% AY2024 力学 第2問 転記（問題のみ）
\section*{第2問〔II〕}

図2のように，半径 $r$，慣性モーメント $I$ の滑車を天井に固定し，糸を巻き付け，糸の
一端に質量 $m$ の物体を取り付け静止させたのち，静かに手を離したところ，物体は鉛直
下向きに落下した。このとき，次の問い (1)〜(3) に答えよ。ここで，重力加速度の大きさを
$g$ とし，糸の質量と太さは無視できるものとする。

\begin{enumerate}[label=(\arabic*)]
  \item 物体の加速度を求めよ。
  \item 糸の張力を求めよ。
  \item 最初，物体が静止していた高さを原点とし，原点から鉛直下向き $h$ の位置に物体が
        きたときの滑車の角速度を求めよ。
\end{enumerate}

\begin{center}
% 図2：天井に固定した滑車（半径r,慣性モーメントI）に糸、端に質量mの物体
\begin{tikzpicture}[>=stealth, scale=1.0]
  % 天井
  \draw[line width=1.4pt] (-0.6,3.0) -- (3.6,3.0);
  % 滑車を吊る短い軸
  \draw (2.2,3.0) -- (2.2,2.8);
  % 滑車
  \draw (2.2,1.9) circle (0.9);
  \fill (2.2,1.9) circle (1.2pt);
  \draw[dashed] (2.2,1.9) -- (1.3,1.9); \node[above] at (1.75,1.9) {$r$};
  \node[right] at (3.1,1.5) {$I$};
  % 糸（滑車左側から鉛直に垂れ、物体へ）
  \draw (1.3,1.9) -- (1.3,0.4);
  \draw (0.9,-0.2) rectangle (1.7,0.4); \node[left] at (0.9,0.1) {$m$};
  \node at (3.0,-0.1) {図 2};
\end{tikzpicture}
\end{center}
